\documentclass[A4,12pt,twoside]{book}
\usepackage{amd}
\usepackage{siunitx}
\sisetup{output-decimal-marker={.}, group-separator={\,}, exponent-product=\cdot}

% %--------------------------------------------------------------------------
% %         General Setting
% %--------------------------------------------------------------------------

\graphicspath{{Images/}{../Images/}} %Path of figures
\setkeys{Gin}{width=0.85\textwidth} %Size of figures
\setlength{\cftbeforechapskip}{3pt} %space between items in toc
\setlength{\parindent}{0.5cm} % Idk

\setlist[itemize]{itemsep=5pt}
\setlist[enumerate]{itemsep=5pt}

% Undertittel for kapitler
\newcommand{\chaptersubtitle}[1]{%
    {\centering\large\itshape #1\par}\medskip%
}

\usetikzlibrary{arrows.meta, calc, decorations.pathmorphing}

\usepackage[default]{comicneue}

% Krysshenvisninger til kompendiet (Tabell-, Figur- og Kapittelnumre m.m.)
% Krever at KompendiumTRES0312.aux finnes (dvs. at kompendiet er kompilert minst én gang).
% Alle kompendium-referanser brukes med prefikset "komp-", slik at de aldri kolliderer
% med løsningsforslagets egne \oppgaver-labels (som gjenbruker samme navn som i kompendiet).
\usepackage{xr-hyper}
\externaldocument[komp-]{KompendiumTRES0312}
\renewcommand{\komplabel}[1]{\ref{komp-#1}}

\begin{document}

\newgeometry{top=8cm,bottom=.5in,left=2cm,right=2cm}
\begin{titlepage}
\centering
\vspace*{3cm}
{\Huge\bfseries Løsningsforslag\par}
\vspace{0.5cm}
{\Large TRES0312 --- Fysikk\par}
\vspace{2cm}
{\large Følger som et eget hefte til \textit{Kompendium TRES0312}.\par}
\vspace{0.3cm}
{\normalsize Hver oppgave er gjengitt i sin helhet, med samme ordlyd som i kompendiet (oppgaveteksten er delt mellom de to dokumentene, slik at en endring i kompendiet automatisk følges opp her), etterfulgt av et fullstendig løsningsforslag som følger samme fremgangsmåte som eksemplene i kompendiet.\par}
\end{titlepage}
\restoregeometry

\tableofcontents

% %--------------------------------------------------------------------------
% %         Core of the document
% %--------------------------------------------------------------------------
\subfile{files_losning/0.0.1_HvaErFysikk_LF}
\subfile{files_losning/0.0.2_MatematikkViTrenger_LF}
\subfile{files_losning/0.0.3_Kinematikk_LF}
\subfile{files_losning/0.0.4_Krefter_LF}
\subfile{files_losning/0.0.5_Energi_LF}
\subfile{files_losning/0.0.6_Statikk_LF}

\end{document}
